\documentclass[11pt]{article}
\usepackage{amsmath}
\usepackage{amssymb}
\usepackage{amsthm}
\usepackage{geometry}
\usepackage{booktabs}
\geometry{margin=1in}

\newtheorem{observation}{Observation}
\newtheorem{remark}{Remark}

\title{Thermodynamic Limit and Basin Scaling in the LSR Energy\\
\large Analysis of Hoover et al.\ (ICLR 2025 Workshop)}
\author{}
\date{}

\begin{document}
\maketitle

\section{Setup and Notation}

The LSR (log-sum-ReLU) energy function introduced in~\cite{Hoover2025} is
\begin{equation}\label{eq:lsr}
  E_{\mathrm{LSR}}(\mathbf{x})
  = -\frac{1}{\beta}\log\!\left(
      \epsilon + \sum_{\mu=1}^{M}
      \mathrm{ReLU}\!\left(1 - \frac{\beta}{2}\|\mathbf{x}-\boldsymbol{\xi}_\mu\|^2\right)
    \right),
\end{equation}
where $\boldsymbol{\xi}_\mu \in \{-1,+1\}^D$ are stored memories, $\beta > 0$
is the inverse temperature, and $\epsilon \geq 0$ is a regularization constant.

The support region around memory $\mu$ is
$\mathrm{S}_\mu = \{\mathbf{x} : \|\mathbf{x}-\boldsymbol{\xi}_\mu\| \leq \sqrt{2/\beta}\}$.
Outside $\bigcup_\mu \mathrm{S}_\mu$ (with $\epsilon = 0$) the energy is infinite
and the gradient vanishes.

For comparison, the standard LSE (log-sum-exp) energy is
\begin{equation}\label{eq:lse}
  E_{\mathrm{LSE}}(\mathbf{x})
  = -\frac{1}{\beta}\log\sum_{\mu=1}^{M}
    \exp\!\left(-\frac{\beta}{2}\|\mathbf{x}-\boldsymbol{\xi}_\mu\|^2\right).
\end{equation}

%----------------------------------------------------------------------
\section{Recap of the Main Theorems and Their Intuition}

We first summarize the three theorems of~\cite{Hoover2025} and explain what
they say in plain terms.

\subsection{Theorem~1: Exact retrieval from basins}

\textbf{Statement (paraphrased).}
Let $r = \min_{\mu\neq\nu}\|\boldsymbol{\xi}_\mu - \boldsymbol{\xi}_\nu\|$
be the minimum distance between any two memories.  For basin radius
$\Delta \in (0,r)$ and $\beta = 2/(r-\Delta)^2$, every input
$\mathbf{x}\in S_\mu(\Delta) \equiv \{\mathbf{x}:\|\mathbf{x}-\boldsymbol{\xi}_\mu\|\leq\Delta\}$
is retrieved exactly as $\boldsymbol{\xi}_\mu$ via energy gradient descent.
With an appropriate learning rate, a \emph{single} gradient step suffices.

\textbf{Intuition.}
The ReLU separation function has finite reach: memory~$\mu$ only influences
points within distance $\sqrt{2/\beta}$ from it.  If $\beta$ is large enough,
this reach is $\sqrt{2/\beta} = r - \Delta < r$.  So any point at distance
$\leq \Delta$ from memory~$\mu$ is within the support of~$\mu$ but
\emph{outside} the support of every other memory (since the nearest other
memory is at distance $\geq r$ away, and $\Delta + (r-\Delta) = r$).
The gradient then points purely toward~$\boldsymbol{\xi}_\mu$, with no
interference from other memories.

\textbf{Hidden constraint: $\Delta \leq r/2$.}
Note a subtlety glossed over in~\cite{Hoover2025}.  The proof claims that for
$\mathbf{x}\in S_\mu(\Delta)$, we have $B(\mathbf{x}) = \{\mu\}$, meaning
$\mathbf{x}$ is in the support of memory~$\mu$.  This requires
$\|\mathbf{x}-\boldsymbol{\xi}_\mu\| \leq \sqrt{2/\beta} = r - \Delta$.
But since $\|\mathbf{x}-\boldsymbol{\xi}_\mu\|\leq\Delta$ by definition of the
basin, we need $\Delta \leq r - \Delta$, i.e.,
\begin{equation}
  \Delta \;\leq\; \frac{r}{2}.
\end{equation}
If $\Delta > r/2$, then points at the outer edge of the ``basin'' are
\emph{outside} the support of their own memory --- they sit in a dead zone
with zero gradient.  So the effective basin can never exceed half the
inter-memory distance.

\subsection{What happens when $\beta \neq 2/(r-\Delta)^2$?}

The prescribed value $\beta = 2/(r-\Delta)^2$ is a knife-edge.  It is
instructive to consider what happens on either side, since the paper does not
discuss this.  The three regimes are controlled by the support radius
$R = \sqrt{2/\beta}$:

\medskip
\textbf{Case 1: $\beta > 2/(r-\Delta)^2$ \;(``too cold'', $R < r-\Delta$).}
The support radius shrinks below $r - \Delta$.  This is \emph{more} than
enough to isolate each memory (neighboring supports are still disjoint), so
there is no interference.  However, the effective basin of attraction also
shrinks: it is no longer $\Delta$ but the support radius $R = \sqrt{2/\beta}$
itself.  Any point farther than $R$ from all memories is in a dead zone.

In this regime:
\begin{itemize}
  \item Retrieval is exact for $\|\mathbf{x}-\boldsymbol{\xi}_\mu\| \leq R$,
  \item but the basin is \emph{smaller} than desired ($R < r - \Delta < \Delta$
        for $\Delta$ near $r/2$),
  \item and a larger fraction of space is dead (zero gradient).
\end{itemize}

\textbf{Case 2: $\beta = 2/(r-\Delta)^2$ \;(the Theorem~1 prescription,
$R = r-\Delta$).}
This is the optimal trade-off: the support radius is as large as possible
while keeping each basin free of interference from other memories.  Every
$\mathbf{x}\in S_\mu(\Delta)$ sees exactly one memory, retrieval is exact.

\textbf{Case 3: $\beta < 2/(r-\Delta)^2$ \;(``too hot'', $R > r-\Delta$).}
The support balls now overlap: a point $\mathbf{x}$ near memory~$\mu$ may
also fall within the support of a neighboring memory~$\nu$ (specifically,
whenever $\|\mathbf{x}-\boldsymbol{\xi}_\nu\| \leq R$ while
$\|\boldsymbol{\xi}_\mu - \boldsymbol{\xi}_\nu\| = r < 2R$).  In the overlap
region, $|B(\mathbf{x})| > 1$, and the gradient becomes an average over
multiple memories.  Gradient descent no longer converges to any individual
$\boldsymbol{\xi}_\mu$ but instead to the centroid
$|B|^{-1}\sum_{\mu\in B}\boldsymbol{\xi}_\mu$ --- an \emph{emergent minimum}.

In this regime:
\begin{itemize}
  \item Original memories are no longer guaranteed to be retrieved from
        their full basin $S_\mu(\Delta)$,
  \item but they remain local minima (at $\mathbf{x} = \boldsymbol{\xi}_\mu$
        the gradient is still zero),
  \item and new local minima (centroids) appear in the overlap regions.
\end{itemize}
This is exactly the regime of Theorem~3.

\medskip
\textbf{Summary: the three regimes.}
\begin{center}
\setlength{\tabcolsep}{4pt}
\begin{tabular}{lccc}
\toprule
& \textbf{Too cold} & \textbf{Knife-edge} & \textbf{Too hot} \\
& $\beta > \frac{2}{(r-\Delta)^2}$ & $\beta = \frac{2}{(r-\Delta)^2}$ &
  $\beta < \frac{2}{(r-\Delta)^2}$ \\
\midrule
Support radius $R$ & $< r - \Delta$ & $= r - \Delta$ & $> r - \Delta$ \\
Supports overlap? & No & No (tangent) & Yes \\
Effective basin & $R < \Delta$ & $\Delta$ & $< \Delta$ (contaminated) \\
Emergent minima & None & None & Yes (centroids) \\
Dead zone fraction & Larger & Moderate & Smaller \\
\bottomrule
\end{tabular}
\end{center}

\begin{remark}[Monotonicity illusion]
One might expect that increasing $\beta$ (``lowering temperature'') always
sharpens retrieval.  In LSE this is true: larger $\beta$ makes the softmax
peakier, which improves basin size.  In LSR the opposite happens at large
$\beta$: the support shrinks and the basin \emph{decreases}, because the
ReLU simply cuts off rather than smoothly suppressing distant memories.
The finite support fundamentally changes the role of temperature.
\end{remark}

\subsection{Theorem~2: Exponential memory capacity}

\textbf{Statement (paraphrased).}
Consider $M$ memories sampled uniformly from $\{-1,+1\}^D$.  With probability
at least $\delta$, all memories are retrievable as per Theorem~1, provided
\begin{equation}
  M = O\!\left(\sqrt{\log(1/\delta)}\;\exp(\alpha^2 D/2)\right),
\end{equation}
where $\alpha \in (0,1)$ is a parameter controlling the trade-off between
capacity and basin size.

\textbf{What is $\alpha$?}
The parameter $\alpha$ is the \emph{maximum cosine similarity} between any
pair of stored memories.  In the proof (Lemma~1 of~\cite{Hoover2025}),
random binary vectors $\boldsymbol{\xi}_\mu, \boldsymbol{\xi}_\nu \in
\{-1,+1\}^D$ satisfy
\begin{equation}
  \frac{|\langle \boldsymbol{\xi}_\mu, \boldsymbol{\xi}_\nu\rangle|}
       {\|\boldsymbol{\xi}_\mu\|\,\|\boldsymbol{\xi}_\nu\|}
  \;\leq\; \alpha
\end{equation}
with high probability (by concentration of measure).  This translates into a
lower bound on the minimum pairwise Euclidean distance:
\begin{equation}\label{eq:r-bound}
  r = \min_{\mu\neq\nu}\|\boldsymbol{\xi}_\mu - \boldsymbol{\xi}_\nu\|
  \;\geq\; \sqrt{2D(1-\alpha)}.
\end{equation}
Smaller $\alpha$ means the memories are more spread out (closer to
orthogonal), giving larger $r$ and larger basins, but fewer memories can be
packed.  Larger $\alpha$ allows more memories at the cost of tighter spacing.
The capacity $M \sim \exp(\alpha^2 D/2)$ grows with $\alpha$: the more
overlap you tolerate, the more patterns you can store.

\textbf{Intuition.}
Random binary vectors in high dimensions are nearly orthogonal: their cosine
similarity concentrates around zero with fluctuations of order $1/\sqrt{D}$.
This means $\sim\!\exp(cD)$ of them can be packed into $\{-1,+1\}^D$ with all
pairwise cosines below~$\alpha$.  Since they are well-separated, Theorem~1
applies to each.

\subsection{Theorem~3: Emergent local minima}

\textbf{Statement (paraphrased).}
Under the same configuration, for any $\mathbf{x}$ that lies in the support
of more than one memory ($|B(\mathbf{x})| > 1$ where
$B(\mathbf{x}) = \{\mu : \|\mathbf{x}-\boldsymbol{\xi}_\mu\|\leq\sqrt{2/\beta}\}$),
the energy gradient descent converges to the centroid
$\frac{1}{|B(\mathbf{x})|}\sum_{\mu\in B(\mathbf{x})}\boldsymbol{\xi}_\mu$,
which is \emph{not} any stored memory.

\textbf{Intuition.}
In the overlap region where two or more support balls intersect, all the
corresponding memories pull on the state simultaneously.  Since the ReLU
separation function is linear (not exponential), no single memory dominates ---
the gradient is a balanced average.  The fixed point is therefore the centroid
of the contributing memories: a \emph{new pattern} that was never stored.
The total number of such emergent minima (one per non-empty subset of memories
with intersecting supports) can be as large as $2^M$.

This is impossible in LSE: the exponential separation function always makes
one memory dominate the softmax, suppressing the others.

%----------------------------------------------------------------------
\section{Implicit Scaling: $\beta$ Cannot Be Fixed as $D\to\infty$}

\subsection{What the theorems require}

From Theorem~2, the minimum pairwise distance satisfies
$r \geq \sqrt{2D(1-\alpha)}$, and exact retrieval with basin radius
$\Delta \in (0,r)$ requires (Theorem~1)
\begin{equation}\label{eq:beta-bound}
  \beta \;=\; \frac{2}{(r-\Delta)^2}.
\end{equation}

\subsection{Why $\Delta$ must be proportional to $r$}

The basin radius $\Delta$ is the maximum corruption from which a memory can be
recovered.  To be a useful associative memory, $\Delta$ must scale with the
natural noise level.  Consider a stored binary pattern
$\boldsymbol{\xi}_\mu\in\{-1,+1\}^D$ corrupted by flipping a fraction~$f$ of
its bits.  The corrupted vector $\tilde{\boldsymbol{\xi}}$ satisfies
\begin{equation}
  \|\tilde{\boldsymbol{\xi}} - \boldsymbol{\xi}_\mu\|^2
  = 4fD
  \quad\Longrightarrow\quad
  \|\tilde{\boldsymbol{\xi}} - \boldsymbol{\xi}_\mu\| = 2\sqrt{fD}.
\end{equation}
Since $r \sim \sqrt{2D}$, this corruption distance is
$\sim (2\sqrt{f}/\sqrt{2})\,r = \sqrt{2f}\;r$.
Retrieving a pattern with fraction~$f$ of bits flipped requires
$\Delta \geq 2\sqrt{fD} \propto r$.  If $\Delta$ did not grow with $r$ (and
hence with $\sqrt{D}$), the memory could only correct a \emph{vanishing}
fraction of bit flips as $D\to\infty$ --- an increasingly useless memory.

More directly: the basin must be large enough that no other memory falls inside
it, which requires $\Delta < r$, but small enough that the support ball
$\sqrt{2/\beta} = r - \Delta$ stays within the exclusive zone.  The only
consistent choice is $\Delta = \eta\, r$ for some fixed $\eta \in (0,1)$.

\subsection{$\beta = O(1/D)$ follows immediately}

With $\Delta = \eta\, r$ and $r \sim \sqrt{2D(1-\alpha)}$:
\begin{equation}
  \beta
  \;=\; \frac{2}{(r - \eta r)^2}
  \;=\; \frac{2}{(1-\eta)^2 r^2}
  \;\sim\; \frac{2}{(1-\eta)^2 \cdot 2D(1-\alpha)}
  \;=\; \frac{1}{(1-\eta)^2(1-\alpha)\,D}.
\end{equation}
More generally, for any fixed ratio $\Delta/r < 1$, we obtain
$\beta = O(1/D)$.  The inverse temperature must vanish as $D\to\infty$.

This is never stated explicitly in~\cite{Hoover2025}, but it is embedded in
every theorem.

\subsection{Consequence for the support radius}

The support radius of each memory is
\begin{equation}
  R_{\mathrm{supp}} = \sqrt{2/\beta} = O(\sqrt{D}),
\end{equation}
which matches the natural length scale of the hypercube $\{-1,+1\}^D$
(where all points lie at distance $\sqrt{D}$ from the origin and typical
inter-point distances are $\sim\!\sqrt{2D}$).

If instead $\beta$ were held fixed as $D\to\infty$:
\begin{itemize}
  \item $R_{\mathrm{supp}} = \sqrt{2/\beta} = O(1)$, a fixed constant;
  \item inter-memory distance $r \sim \sqrt{2D} \to \infty$;
  \item each memory becomes an isolated island of finite energy in a space
        where all scales grow as $\sqrt{D}$;
  \item the fraction of space with any gradient signal goes to zero
        exponentially.
\end{itemize}

\subsection{Contrast with LSE}

For the LSE energy~\eqref{eq:lse}, $\beta$ can remain $O(1)$ as $D\to\infty$.
The exponential separation function has infinite support: every point
$\mathbf{x}\in\mathbb{R}^D$ receives a nonzero gradient pointing toward the
softmax-weighted centroid of the memories.  The basins of attraction have
radius $O(\sqrt{D})$ even with $\beta = O(1)$.  No dead zones exist.

%----------------------------------------------------------------------
\section{The Thermodynamic Limit}

\subsection{Mean-field scaling}

Define a rescaled inverse temperature
\begin{equation}
  \tilde\beta = \beta D, \qquad \text{so that } \beta = \tilde\beta/D,
\end{equation}
and hold $\tilde\beta$ fixed as $D\to\infty$.  The energy becomes
\begin{equation}
  E_{\mathrm{LSR}}(\mathbf{x})
  = -\frac{D}{\tilde\beta}\,\log\!\left(
      \epsilon + \sum_{\mu=1}^{M}
      \mathrm{ReLU}\!\left(
        1 - \frac{\tilde\beta}{2D}\|\mathbf{x}-\boldsymbol{\xi}_\mu\|^2
      \right)
    \right).
\end{equation}

Since $\|\mathbf{x}-\boldsymbol{\xi}_\mu\|^2 = O(D)$ for typical
configurations, the ReLU argument is $1 - O(\tilde\beta) = O(1)$, and
\begin{equation}
  E_{\mathrm{LSR}} = -\frac{D}{\tilde\beta}\,\log\bigl(O(1)\bigr) = O(D).
\end{equation}

\begin{observation}[Extensivity]
The energy is extensive (proportional to system size $D$), as required for a
well-defined thermodynamic limit.  The correct limit is
\[
  D \to \infty, \qquad \tilde\beta = \beta D \;\text{fixed}, \qquad
  M \sim e^{cD} \;\text{for some } c > 0.
\]
This is analogous to the mean-field scaling $J \sim 1/\sqrt{N}$ in classical
spin glass models.
\end{observation}

\subsection{Comparison of scaling regimes}

\begin{table}[h]
\centering
\begin{tabular}{lcc}
\toprule
\textbf{Quantity} & \textbf{LSE} & \textbf{LSR} \\
\midrule
Inverse temperature $\beta$ & $O(1)$ & $O(1/D)$ \\
Rescaled temperature $\tilde\beta = \beta D$ & grows as $D$ & $O(1)$ \\
Support radius $\sqrt{2/\beta}$ & $\infty$ (global) & $O(\sqrt{D})$ \\
Basin radius $\Delta$ & $O(\sqrt{D})$ & $O(\sqrt{D})$ \\
Memory capacity $M$ & $\sim e^{cD}$ & $\sim e^{cD}$ \\
Energy $E$ & $O(D)$ & $O(D)$ \\
Dead zones & None & Yes \\
\bottomrule
\end{tabular}
\caption{Scaling comparison between LSE and LSR in the $D\to\infty$ limit.}
\label{tab:comparison}
\end{table}

%----------------------------------------------------------------------
\section{The Basin Shrinkage Problem}

\subsection{Absolute vs.\ relative basin size}

In absolute terms, the basin radius $\Delta \sim \sqrt{D}$ \emph{grows} with
dimension.  The support radius $\sqrt{2/\beta}\sim\sqrt{D}$ also grows.
All natural length scales ($r$, $\Delta$, $R_{\mathrm{supp}}$, $\|\boldsymbol{\xi}_\mu\|$)
are $O(\sqrt{D})$, so the geometry is self-similar across dimensions.

In relative terms, however, the \emph{volume fraction} of a $D$-dimensional
ball of radius $c\sqrt{D}$ inside the hypercube $[-1,1]^D$ behaves as
\begin{equation}
  \frac{\mathrm{Vol}(B_{c\sqrt{D}})}{\mathrm{Vol}([-1,1]^D)}
  \;\sim\;
  \left(\frac{\pi e\, c^2}{4}\right)^{D/2}
  \!\!\cdot\; \text{(subexponential prefactor)}.
\end{equation}
For moderate $c$ this ratio decays exponentially with $D$.  But this is
equally true for LSE --- the basins of attraction for LSE are also balls of
radius $O(\sqrt{D})$ covering an exponentially small fraction of the space.

\subsection{The qualitative difference: dead zones}

The fundamental distinction is not basin size but what happens \emph{outside}
all basins:

\begin{center}
\begin{tabular}{lcc}
\toprule
& \textbf{LSE} & \textbf{LSR} \\
\midrule
Outside all basins & Gradient exists & \textbf{Gradient is zero} \\
Random initialization & Always converges & May land in dead zone \\
Retrieval from corruption & Works if $\|\delta\| < \Delta$ &
  Same condition \\
\bottomrule
\end{tabular}
\end{center}

For the \textbf{retrieval task} (starting from a corrupted version of a
stored memory with bounded noise), both models work identically well:
the corrupted input is at distance $O(\sqrt{D})$ from the true memory,
which lies within the basin.

For \textbf{random initialization} or unconstrained generation, LSR has a
fundamental problem: any initial point outside $\bigcup_\mu \mathrm{S}_\mu$
has zero gradient and will not move under gradient descent.  The model is
blind in these regions.

\subsection{Emergent minima in the thermodynamic limit}

With $\beta = \tilde\beta/D$ and $\tilde\beta$ fixed, the support radius
$R_{\mathrm{supp}} = \sqrt{2D/\tilde\beta}$ and the inter-memory distance
$r \sim \sqrt{2D}$.  Two memories at distance $r$ have overlapping support
regions whenever $2R_{\mathrm{supp}} > r$, i.e.,
\begin{equation}
  2\sqrt{2D/\tilde\beta} > \sqrt{2D}
  \quad\Longleftrightarrow\quad
  \tilde\beta < 4.
\end{equation}
So for $\tilde\beta < 4$, supports overlap, emergent minima (centroids of
memory subsets) appear, and their number can be as large as $2^M$.  The
basin radius for exact retrieval of the original memories requires
$\tilde\beta > 2/(1 - \Delta/r)^2 \cdot (1/D) \cdot D = 2D/(r-\Delta)^2$;
substituting $r \sim \sqrt{2D}$ and $\Delta = \eta r$ for $\eta < 1$:
\begin{equation}
  \tilde\beta \;\geq\; \frac{1}{(1-\eta)^2(1-\alpha)}.
\end{equation}
Thus the regime of simultaneous memorization and generation is
\begin{equation}
  \frac{1}{(1-\eta)^2(1-\alpha)} \;\leq\; \tilde\beta \;<\; 4,
\end{equation}
which exists for sufficiently small $\eta$ and $\alpha$.

%----------------------------------------------------------------------
\section{Summary}

\begin{enumerate}
  \item \textbf{$\beta$ is not fixed.}
    The theorems in~\cite{Hoover2025} implicitly require $\beta = O(1/D)$.
    At fixed $\beta$, the support regions become negligibly small relative
    to the ambient space as $D\to\infty$.

  \item \textbf{A thermodynamic limit exists} with mean-field scaling
    $\tilde\beta = \beta D$ held fixed.  The energy is extensive, $E = O(D)$,
    and memory capacity remains exponential.

  \item \textbf{Basins grow as $\sqrt{D}$} in absolute terms, matching all
    other length scales.  The exponential shrinkage of volume fractions
    affects both LSE and LSR equally.

  \item \textbf{Dead zones are the real cost} of finite support.
    Unlike LSE, which provides a gradient signal everywhere, LSR has regions
    of zero gradient that cover an increasing fraction of the space.
    This is the price paid for the emergent minima property.

  \item \textbf{Finite support is a double-edged sword.}
    The same mechanism that creates emergent local minima (overlapping
    finite support regions) also creates dead zones (regions outside all
    supports).  The paper's strength and weakness are two sides of the same
    coin.
\end{enumerate}

\begin{thebibliography}{9}
\bibitem{Hoover2025}
  B.~Hoover, K.~Balasubramanian, D.~Krotov, and P.~Ram,
  ``Dense Associative Memory with Epanechnikov Energy,''
  in \emph{New Frontiers in Associative Memory} (Workshop at ICLR 2025),
  2025.
  Available at: \texttt{https://openreview.net/forum?id=LSR\_DenseAM}

\bibitem{Krotov2016}
  D.~Krotov and J.~J.~Hopfield,
  ``Dense associative memory for pattern recognition,''
  \emph{Advances in Neural Information Processing Systems}, vol.~29, 2016.

\bibitem{Ramsauer2021}
  H.~Ramsauer, B.~Sch\"afl, J.~Lehner, P.~Seidl, M.~Widrich, L.~Gruber,
  et al.,
  ``Hopfield networks is all you need,''
  in \emph{International Conference on Learning Representations}, 2021.

\bibitem{Lucibello2024}
  C.~Lucibello and M.~M\'ezard,
  ``Exponential capacity of dense associative memories,''
  \emph{Physical Review Letters}, vol.~132, no.~7, p.~077301, 2024.

\bibitem{Epanechnikov1969}
  V.~A.~Epanechnikov,
  ``Non-parametric estimation of a multivariate probability density,''
  \emph{Theory of Probability \& Its Applications}, vol.~14, no.~1,
  pp.~153--158, 1969.

\bibitem{WandJones1994}
  M.~P.~Wand and M.~C.~Jones,
  \emph{Kernel Smoothing}, CRC Press, 1994.
\end{thebibliography}

\end{document}
